\documentclass[11pt,a4paper]{article}

\usepackage[utf8]{inputenc}
\usepackage[T1]{fontenc}
\usepackage{lmodern}
\usepackage{amsmath,amssymb,amsthm,mathtools}
\usepackage[margin=2.5cm]{geometry}
\usepackage{hyperref}
\usepackage{cleveref}
\usepackage{enumitem}
\usepackage{booktabs}
\usepackage{xcolor}
\usepackage{fancyhdr}

\pagestyle{fancy}
\fancyhf{}
\fancyhead[L]{\small\textsc{SCT Theory --- NT-3 Literature}}
\fancyhead[R]{\small\thepage}
\renewcommand{\headrulewidth}{0.4pt}

\hypersetup{
  colorlinks=true,
  linkcolor=blue!60!black,
  citecolor=green!40!black,
  urlcolor=blue!60!black,
  pdftitle={NT-3: Spectral Dimension Literature Review},
  pdfauthor={David Alfyorov},
}

\title{%
  \textbf{NT-3: Spectral Dimension in Quantum Gravity}\\[4pt]
  \large Literature Review and Conventions
}
\author{David Alfyorov}
\date{March 2026}

\begin{document}
\maketitle

\begin{center}
and workflow support.}
\end{center}

\tableofcontents

\section{Overview}

The spectral dimension $d_S(\sigma)$ quantifies how spacetime dimensionality
changes as a function of the diffusion scale $\sigma$. This document reviews
conventions, key formulas, and benchmarks from 20 papers spanning CDT,
Asymptotic Safety, Horava--Lifshitz gravity, Stelle gravity, LQG,
causal sets, and noncommutative geometry.

\section{Definition}

The spectral dimension is defined via the return probability $P(\sigma)$
of a fictitious diffusion process:
\begin{equation}\label{eq:ds-def}
  d_S(\sigma) = -2\,\frac{\mathrm{d}\ln P(\sigma)}{\mathrm{d}\ln\sigma}\,.
\end{equation}
For a standard $d$-dimensional flat manifold with Laplacian $\Delta$,
the heat kernel trace gives $P(\sigma) = (4\pi\sigma)^{-d/2}$, yielding
$d_S = d$.

\section{Multiple Definitions in Modified Gravity}

Three distinct definitions arise in theories with a modified propagator
$G(k^2)$:
\begin{itemize}
  \item \textbf{Heat kernel:} $P_{\mathrm{HK}}(\sigma) = \int\frac{\mathrm{d}^dk}{(2\pi)^d}\,
    e^{-\sigma k^2}$ gives $d_S = d$ (unmodified).
  \item \textbf{SVW:} $P_{\mathrm{SVW}}(\sigma) = \int\frac{\mathrm{d}^dk}{(2\pi)^d}\,
    e^{-\sigma\omega^2(k)}$ where $\omega^2(k) = k^2\Pi_{\mathrm{TT}}(k^2/\Lambda^2)$.
    For $\omega^2 \sim k^{2\alpha}$: $d_S = d/\alpha$ (Sotiriou--Visser--Weinfurtner, 1105.6098).
  \item \textbf{CMN (propagator-weighted):} $P_{\mathrm{CMN}}(\sigma) = \int\frac{\mathrm{d}^dk}{(2\pi)^d}\,
    G(k^2)\,e^{-\sigma k^2}$. For $G \sim 1/k^{2n}$: $d_S = d - 2n$
    (Calcagni--Modesto--Nardelli, 1408.0199).
\end{itemize}

\section{Benchmarks}

\begin{table}[h]
\centering
\begin{tabular}{llccl}
\toprule
Approach & $d_S(\mathrm{IR})$ & $d_S(\mathrm{UV})$ & Method & Reference \\
\midrule
GR & 4 & 4 & Exact & --- \\
CDT & $4.02 \pm 0.10$ & $1.80 \pm 0.25$ & Monte Carlo & hep-th/0505113 \\
Asymptotic Safety & 4 & 2 & Exact (NGFP) & hep-th/0511260 \\
Horava--Lifshitz ($z=3$) & 4 & 2 & Exact & 0902.3657 \\
Stelle gravity & 4 & 2 & Exact & 1408.0199 \\
LQG (polymer) & 4 & 2 & Semi-analytic & 0912.0220 \\
Causal sets & 4 & $>4$ & Numerical & 1311.2530 \\
Spectral action (ASZ, full) & 4 & 0 & Analytic & 1410.7999 \\
\bottomrule
\end{tabular}
\caption{Spectral dimension values in quantum gravity approaches.}
\label{tab:benchmarks}
\end{table}

\section{SCT-Specific Input}

From the verified NT-4a results, the SCT propagator has
$\Pi_{\mathrm{TT}}(z) = 1 + \tfrac{13}{60}\,z\,\hat{F}_1(z)$
with $\Pi_{\mathrm{TT}}(0) = 1$ and $\Pi_{\mathrm{TT}} \to -83/6$
as $z \to +\infty$. The ghost zero lies at $z_0 \approx 2.4148$.

The \emph{saturation} of $\Pi_{\mathrm{TT}}$ at a constant means the
propagator scales as $1/k^2$ in the UV, unlike Stelle gravity where
$G \sim 1/k^4$. This is the key structural difference that determines
the SCT spectral dimension.

\section{Bibliography}

20 papers reviewed; key references: Ambjorn--Jurkiewicz--Loll (hep-th/0505113),
Lauscher--Reuter (hep-th/0511260), Horava (0902.3657),
Sotiriou--Visser--Weinfurtner (1105.6098),
Calcagni--Modesto--Nardelli (1408.0199),
Alkofer--Saueressig--Zanusso (1410.7999),
Modesto (1209.5129).

\end{document}
